% Source: https://artofproblemsolving.com/community/c3h1488481_really_confusing_problem
% Stable identity: aops-topic-1488481-q01
Let the population at $t=0$ be $P$. An increase of $i\%$ multiplies it by $1+\dfrac{i}{100}$, so at $t=1$ it is $P\left(1+\dfrac{i}{100}\right)$. The second increase of $j\%$ is a percentage of this new population, so at $t=2$ the population is
\[
P\left(1+\dfrac{i}{100}\right)\left(1+\dfrac{j}{100}\right)=P\left(1+\dfrac{i+j}{100}+\dfrac{ij}{10000}\right)=P\left(1+\dfrac{1}{100}\left(i+j+\dfrac{ij}{100}\right)\right).
\]
So the total increase is $\left(i+j+\dfrac{ij}{100}\right)\%$, answer $\boxed{\textbf{(D)}}$. (For example, $10\%$ then $20\%$ gives $1.1\cdot1.2=1.32$, an increase of $32\%=10+20+\dfrac{200}{100}$ percent.)
